% titel: Blatt 9, A2: Integrationsregeln
% Teil II: der durchgerechnete Weg.

\subsection*{a) Partielle Integration}

\textbf{$\int_0^\pi (x+1)^2\cos x\,\dd x$:} Für das Polynom $f(x)=(x+1)^2$
und $g'(x)=\cos x$ ist $g(x)=\sin x$, also
\begin{align}
  \int_0^\pi (x+1)^2\cos x\,\dd x
  &= \Bigl[(x+1)^2\sin x\Bigr]_0^\pi
     - 2\int_0^\pi (x+1)\,\sin x\,\dd x .
\end{align}
Randterm: $(x+1)^2\sin x\big|_0^\pi = (\pi+1)^2\cdot0-1\cdot0=0$. Nochmals
partiell integrieren mit $f(x)=x+1$, $g'=\sin x$:
\begin{align}
  \int (x+1)\sin x\,\dd x
  &= -(x+1)\cos x + \int \cos x\,\dd x
   = -(x+1)\cos x + \sin x .
\end{align}
Eingesetzt:
\begin{align}
  \int_0^\pi (x+1)^2\cos x\,\dd x
  &= -2\Bigl[\overbrace{-(x+1)\cos x+\sin x}\Bigr]_0^\pi \\
  &= -2\Bigl[\bigl(-(\pi+1)(-1)+0\bigr)-\bigl(-1\cdot1+0\bigr)\Bigr]
   = -2\bigl[(\pi+1)-(-1)\bigr]
   = -2(\pi+2).
\end{align}

\textbf{$\int_1^4 \frac{\ln x}{\sqrt x}\,\dd x$:} Wähle $f=\ln x$, $g'=x^{-1/2}$
$\Rightarrow$ $g=2x^{1/2}$:
\begin{align}
  \int \frac{\ln x}{\sqrt x}\,\dd x
  &= 2\sqrt x\,\ln x - \int 2x^{1/2}\cdot\frac1x\,\dd x
   = 2\sqrt x\,\ln x - \int 2x^{-1/2}\,\dd x \\
  &= 2\sqrt x\,\ln x - 4\sqrt x .
\end{align}
Mit den Grenzen:
\begin{align}
  \Bigl[2\sqrt x\,\ln x - 4\sqrt x\Bigr]_1^4
  &= (4\ln4 - 8) - (0-4) = 4\ln4 - 4 = 4(\ln4-1) = 8\ln2 - 4 .
\end{align}

\begin{ergebnis}
$\displaystyle\int_0^\pi (x+1)^2\cos x\,\dd x = -2(\pi+2),\qquad
\int_1^4\frac{\ln x}{\sqrt x}\,\dd x = 8\ln2-4.$
\end{ergebnis}

\subsection*{b) Substitution}

\textbf{$\int_0^1 x\sqrt{1+x^2}\,\dd x$:} $u=1+x^2$, $\dd u=2x\,\dd x$
$\Rightarrow$ $x\,\dd x=\frac12\,\dd u$; Grenzen
$x{=}0{\to}u{=}1$, $x{=}1{\to}u{=}2$:
\begin{align}
  \int_0^1 x\sqrt{1+x^2}\,\dd x
  = \frac12\int_1^2 u^{1/2}\,\dd u
  = \frac12\Bigl[\frac23 u^{3/2}\Bigr]_1^2
  = \frac13\bigl(2\sqrt2-1\bigr) .
\end{align}

\textbf{$\int_0^\pi \cos^3 x\,\sin x\,\dd x$:} $u=\cos x$, $\dd u=-\sin x\,\dd x$,
Grenzen $0{\to}1$, $\pi{\to}{-}1$:
\begin{align}
  \int_0^\pi \cos^3 x\,\sin x\,\dd x
  = \int_{1}^{-1} u^3(-\dd u)
  = \int_{-1}^{1} u^3\,\dd u = 0 ,
\end{align}
denn $u^3$ ist ungerade (s. Blatt 9, A5), das Integral über $[-1,1]$ ist null.

\textbf{$\int_1^e \frac{(\ln x)^2}{x}\,\dd x$:} $u=\ln x$, $\dd u=\frac1x\,\dd x$,
Grenzen $1{\to}0$, $e{\to}1$:
\begin{align}
  \int_1^e \frac{(\ln x)^2}{x}\,\dd x
  = \int_0^1 u^2\,\dd u = \frac{u^3}{3}\Big|_0^1 = \frac13 .
\end{align}

\begin{ergebnis}
$\displaystyle\int_0^1 x\sqrt{1+x^2}\,\dd x=\frac13(2\sqrt2-1),\qquad
\int_0^\pi \cos^3 x\,\sin x\,\dd x=0,\qquad
\int_1^e\frac{(\ln x)^2}{x}\,\dd x=\frac13.$
\end{ergebnis}

\begin{kontrolle}
\textbf{Grössenordnung:} $\frac13(2\sqrt2-1)\approx0.609$ passt zu einem
positiven Integranden $x\sqrt{1+x^2}\in[0,\sqrt2]$. \textbf{Null.} $\cos^3x\sin x$ ist auf $[0,\pi]$
symmetrisch mit Vorzeichenwechsel um $x=\frac\pi2$, das Integral ist $0$ —
konsistent. \textbf{Kleinheit:} $(\ln x)^2/x$ positiv, kleines Ergebnis
$\frac13$ plausibel.
\end{kontrolle}

\begin{merksatz}
Partielle Integration, wenn ein Faktor beim Ableiten einfacher wird
(Polynom$\cdot\sin/\cos$); Substitution, wenn $\varphi'(x)$ als Faktor schon
da ist ($x\dd x$ zu $x^2$, $\frac1x\dd x$ zu $\ln x$, $\sin\dd x$ zu $\cos$).
Grenzen immer mit umrechnen.
\end{merksatz}
